\documentclass[9pt,a4paper]{article}

\usepackage[T1]{fontenc}
\usepackage[utf8]{inputenc}
\usepackage[cyr]{aeguill}
\usepackage[francais]{babel}

% %\renewcommand{\baselinestretch}{1.5}  % line spacing : 1,5
% \setlength{\oddsidemargin}{2,5cm}	% Marge gauche sur pages impaires
% \setlength{\evensidemargin}{2,5cm}	% Marge gauche sur pages paires
% \setlength{\textwidth}{16cm}	% Largeur de la zone de texte (17cm)
% \setlength{\topmargin}{2,5cm}	% Pas de marge en haut

% Les packages
%\usepackage[french]{babel}
\usepackage{fancybox}
\usepackage{graphics}
\usepackage{color}
\usepackage{latexsym}
\usepackage{amsmath,amsthm,amssymb}
\usepackage[plain]{fullpage} 
\usepackage{verbatim}
\usepackage{times,epsfig}
\usepackage{listings}
\usepackage{multirow}
\usepackage{amsmath,amsthm,amssymb,amscd,txfonts,bbm} % RAJOUT ICI !!!
\usepackage{graphics,epsfig} % RAJOUT ICI !!!
%\usepackage{algorithmic}
\usepackage[final]{pdfpages} 
\usepackage{fancyvrb}
\usepackage{enumerate}
\usepackage{listingsutf8}
\usepackage{tikz}
\usetikzlibrary{arrows,shapes,snakes,automata,backgrounds,petri,positioning}
\usepackage{tikz-qtree}
\usepackage{comment}


\newenvironment{qv}
{\quote\Verbatim}
{\endquote\endVerbatim}

\usepackage{lastpage}

\newcommand{\terme}[1]{\texttt{#1}}
\newcommand{\termSet}[1]{\{\texttt{#1}\}}
\newcommand{\guill}[1]{<<\,#1\,>>}
% Macro pour l'inclusion de figure  
\newcommand{\fig}[2]{%
  \begin{figure*}[hbt]
   \begin{center}
      \includegraphics[width=6cm]{#1}
  \caption{\label{#1}  #2}
  \end{center}
 \end{figure*}
}
\newcommand{\figTex}[2]{%
  \begin{figure*}[hbt]
 \centerline{\input{#1.pstex_t}}
  \caption{\label{#1}  #2}
 \end{figure*}
}

% Macro commentaires
\definecolor{turquoise}{rgb}{0.20 ,0.60, 0.60}
\definecolor{vert}{rgb}{0.40,  0.60 ,0.00}
\definecolor{violet}{rgb}{0.80 , 0.00 , 0.80}
\definecolor{bleu}{rgb}{0.20,  0.20,  0.80}
\definecolor{rose}{rgb}{1.00,  0.20,  0.40}
\def\remVG#1{\textcolor{vert}{\texttt{Virginie: #1}}}
\def\comVG#1{\begin{footnotesize}\textcolor{vert}{\texttt{Precision de Virginie: #1}}\end{footnotesize}}

\lstset{basicstyle=\footnotesize,showstringspaces=false,language=XML,breaklines=false,columns=fullflexible,frame=shadowbox, captionpos=b}

\newenvironment{solution}[0]{
~\\ \color{red} \textbf{Solution :}\color{black}}{%
}
%\excludecomment{solution}


\parindent=0pt           
                             
\usepackage{url}

\newtheorem{Exercice}{Definition}

\lstset{language=XML,basicstyle=\footnotesize,
numberstyle=\footnotesize, 
breaklines=true
showspaces=false,         
showstringspaces=false,
showtabs=false,
frame=single,
tabsize=4}   


\begin{document}

\title{Examen NFE204 - Session 1}
\date{29 janvier 2013}

\maketitle

\begin{center}
{\bf \Large Consignes}\\
\begin{tabular}{| p{15cm} |}
\hline
\\
Dur\'ee de l'examen : 2h30\\
Documents non autoris\'es\\
Calculatrice autoris\'ee\\
% Il aurait fallu les pr�venir avant
%\textbf{Les transparents imprim\'es du cours sont autoris\'es. (\textsc{VT :
%mieux pour XML en raison de la syntaxe compliqu\'ee}, \`a discuter)}\\
%Tous les documents au format papier sont autoris\'es. \\ 
Les exercices sont ind\'ependants les uns des autres. Merci de soigner votre \'ecriture.\\
Ce sujet contient \pageref{LastPage} pages.\\


\textit{Important}\,: nous n'évaluons pas les réponses par la syntaxe. ne vous
inquiétez pas d'utiliser un point-virgule à la place dun point d'exclamation ou
autres détails mineurs
\\
\hline
\end{tabular}
\end{center}

\pagebreak

\section*{Première partie : XML}
. 
\subsection*{\Large Exercice 1 - XPath, XQuery, XSLT ({\it 10 points})}

Voici le contenu du fichier \textit{mylocation.xml} :
\vspace{-0.05cm} 
\lstinputlisting[inputencoding=utf8/latin1,basicstyle=\small,language=XML,breaklines=true, caption={Fichier mylocation.xml}]{my_location.xml}

\subsubsection*{XPath}
\begin{enumerate}
\item Quel est le r\'esultat de l'\'evaluation de chacune de ces requêtes XPath sur le document XML ci-dessus ?
\begin{enumerate}
\vspace{-0.1cm}
\item \verb7/agence/studio/description7
\begin{solution}
\begin{verbatim}
<description>beau studio de 20 m ENVIRON a louer vide au 3eme etage d’un immeuble ancien refait neuf lumineux calme emplacement de qualite !</description>
\end{verbatim}
\end{solution}
\vspace{-0.1cm}
\item \verb7/agence/appart[@prix<1000 && @superficie>50]/quartier/text()7
\begin{solution}
$\oslash$
\end{solution}
\end{enumerate}

\item Donner les requêtes XPath permettant de récupérer 
\begin{enumerate}
\vspace{-0.1cm}
\item Les descriptions des studios de plus de 18m$^2$ dans le quartier des paquerettes. 
\begin{solution}
\verb7/agence/studio[@superficie>18][quartier="Les paquerettes"]/description7
\end{solution}
\vspace{-0.1cm}
\item Les appartements de plus de 50m$^2$ ayant au moins chacun un de ces critères : un ascenseur et un parking.
\vspace{-0.1cm}
\begin{solution}
\verb7/agence/appartement[@superficie>50][ascenseur][parking]7
\end{solution}
% \item Le prix des appartements de plus de 2 pièces à vivre
% \begin{solution}
% \verb7/agence/appartement[count(piece) > 2]/@prix7
% \end{solution}
\end{enumerate}
\end{enumerate}

% \subsubsection*{XQuery}
% \begin{enumerate}
%   \item[3.] Donner une requête XQuery utilisant la syntaxe FLOWR
% permettant de r\'ecup\'erer pour chaque location :
% 	\begin{itemize}
% 	  \item En attribut : son prix, sa superficie, le nombre de pièces à vivre, le nombre de pièces d'eau, la ville
% 	  \item En élément fils : la description
% 	\end{itemize}
% \begin{solution}
% \begin{verbatim}
% for $i in document("mylocation.xml")/agence/*
% return
%     <location prix="{$i/@prix}" supercifie="{$i/@superficie}" nbpiece="{count($i/piece)}"
%               ville="{$i/quartier/@ville}">
%     	{$i/description}
%     </location>
% \end{verbatim}
% \end{solution}
% \item[]
% 	\item[4.] Rajouter dans la requête précédente un ordonnencement des locations par prix au m$^2$ croissant ('prix/superficie').
% \begin{solution}
% \begin{verbatim}
% order by ($i/@prix / $i/@superficie)
% \end{verbatim}
% \end{solution}
% \end{enumerate}

\subsubsection*{XSLT}
\begin{enumerate}
	\item[5.] Quel est le code XHML issu de la transformation du fichier
	\textit{mylocation.xml} avec la feuille de style XSLT suivante :

\begin{center}
\begin{minipage}{\textwidth}
\lstinputlisting[inputencoding=utf8/latin1,basicstyle=\small,
language=XSLT,breaklines=true, caption={Feuille de style XSLT}]{feuille.xsl}
\end{minipage} 
\end{center}

\begin{solution}
\begin{verbatim}
<html>
  <body>
    <h1>Liste d Appartements</h1>
    <ul>
    </ul>
    <h1>Liste avec description</h1>
    <ul>
      <li>Appartement de 60 m2, 1200 euros, a Chaville.<br/>
        Salon de 26 m2, Chambre de 10 m2, Chambre de 10 m2, Cuisine equipee de 4 m2, Salle de Bain de 5 m2, toilettes de 2 m2, ascenseur, cave,
        <p><b>Description :</b> bel appartement de 3 pieces 60 m2 environ a louer meuble au 6 eme etage ascenseur dun immeuble tres securise lumineux parfait etat</p></li>
      <li>Appartement de 50 m2, 900 euros, a Chaville.<br/>
        Cuisine/Salon de 30 m2, Chambre de 10 m2, Salle de Bain de 5 m2, toilettes de 2 m2, cave, parking, 
        <p><b>Description :</b> tres beau 2 pieces de 50 m2 parfaitement bien meuble et bon etat, proximite de gare Chaville rive Droite</p></li>
    </ul>
  </body>
</html>
\end{verbatim}
\end{solution}
	\item[6.] Donner le code XSL permettant de générer le pattern en mode '$short$' de l'annonce des appartements. Vous mettrez les valeurs de superficie en m2, prix en euros et ville.
\end{enumerate}
\begin{solution}
\begin{lstlisting}[inputencoding=utf8/latin1,basicstyle=\small,
language=XSLT,breaklines=true]
  <xsl:template match="appart" mode="short">
    <li>
    	<xsl:value-of select="@superficie"/> m2,
    	<xsl:value-of select="@prix"/> euros,
    	a <xsl:value-of select="quartier/@ville"/>
    </li>
  </xsl:template>
\end{lstlisting}
\end{solution}
\subsubsection*{DTD}

\begin{enumerate}
  \item[7.] Proposer une DTD qui vérifie le document \textit{mylocation.xml}
\begin{solution}
\begin{verbatim}
<!ELEMENT agence (appart*, studio*)>
<!ELEMENT appart (piece+, cuisine?, piecedeau+, ascenseur?,
          parking?, cave?, description, quartier)>
<!ELEMENT studio (piece, piecedeau?, ascenseur?, description, quartier)>
<!ELEMENT piece (#PCDATA)>
<!ELEMENT piecedeau (#PCDATA)>
<!ELEMENT ascenseur (EMPTY)>
<!ELEMENT cave (EMPTY)>
<!ELEMENT parking (EMPTY)>
<!ELEMENT description (#PCDATA)>
<!ELEMENT quartier (#PCDATA)>
<!ATTLIST appart id ID #REQUIRED, superficie CDATA #REQUIRED,
         prix CDATA #REQUIRED, etage CDATA #REQUIRED> 
<!ATTLIST studio id ID #REQUIRED, superficie CDATA #REQUIRED,
         prix CDATA #REQUIRED, etage CDATA #REQUIRED>
<!ATTLIST piece superficie CDATA #REQUIRED>
<!ATTLIST piecedeau superficie CDATA #REQUIRED>
<!ATTLIST cuisine superficie CDATA #REQUIRED>
<!ATTLIST cave superficie CDATA #REQUIRED>
<!ATTLIST quartier ville CDATA #REQUIRED>
\end{verbatim}
\end{solution}
\end{enumerate}

\section*{Seconde partie : Gestion de donn\'ees à grande \'echelle}

\subsection*{Exercice 2 - Recherche plein texte ({\it 4 points})}

Nous souhaitons indexer le texte (balise description) des locations de notre agence.

\begin{enumerate}
  \item[8.] Cr\'eer les listes inverses pour les mots suivants : ''appartement, etat, immeuble, louer, lumineux''.\\
  Vous y associerez les valeurs de tf et d'idf (aucun traitement sur le texte).
\begin{solution}
\begin{itemize}
  \item doc1 : 21 termes, doc2 : 19 termes, doc3 : 22
  \item appartement (1/3) : doc1 (1/21)
  \item etat (2/3) : doc1 (1/21), doc2 (1/19)
  \item immeuble (2/3) : doc1 (1/21), doc3 (1/22)
  \item louer (2/3) : doc1 (1/21), doc3 (1/21)
  \item lumineux (2/3) : doc1 (1/21), doc3 (1/22)
\end{itemize}
\end{solution}

  \item[9.] Donner la formule de similarité du cosinus en prenant en compte le {\bf tf$_i$}, {\bf idf$_i$}, le poids {\bf w$_{qi}$} des termes de la requête, la norme du document {\bf Nd} et la norme de la requête {\bf Nq}.

  \item[10.] Donner le classement par cosinus des 3 documents avec la requête suivante~: ``appartement louer lumineux$\hat{~}$2''.\\
  Sachant que les normes des documents sont doc1 : 0.0705, doc2 : 0.0841, doc3: 0.0604, et la norme de la requête : 0.6123.
\begin{solution}
\begin{itemize}
  \item doc1 : $\frac{\frac{1}{21} \times \log(\frac{3}{1}) \times \frac{1}{4} + \frac{1}{21} \times \log(\frac{3}{2}) \times \frac{1}{4} + \frac{1}{21} \times \log(\frac{3}{2}) \times \frac{1}{2}}{0.0705 \times 0.6123}  = 0.0277$ 
  \item doc2 : $0$
  \item doc3 : $\frac{\frac{1}{22} \times \log(\frac{3}{2}) \times \frac{1}{4} + \frac{1}{22} \times \log(\frac{3}{2}) \times \frac{1}{2}}{0.0604 \times 0.6123} = 0.0162$
\end{itemize}
\end{solution}
\end{enumerate}
\subsection*{Exercice 3 - Architecture distribuée ({\it 6 points})}

\begin{enumerate}
  \item Expliquer dans quel type d'architecture il peut être nécessaire d'effectuer de la réconciliation de données. Donner un scénario typique.
  \item Qu'apporte le {\it consistent-hashing} pour des systèmes distribués évoluant souvent ? (Ajout/suppression de serveurs)
  \item Qu'appelle-t'on cohérence à terme ({\it eventual consistancy}) dans un système distribué ? Donner un exemple.
\end{enumerate}

% \subsection*{Exercice 3 - Indexation distribu\'ee ({\it 6 points})}
% 
% La figure~\ref{chord-friends} montre un anneau de distribution de type
% \textit{consistent hashing}. Les ronds oranges représentent des serveurs.
% On veut gérer une collection de documents $D$, chacun doté d'un identifiant
% $id$.
% 
% 
%  \fig{chord-friends}{Situation initiale de l'anneau \textit{Consistent hashing}}
%  
% \begin{enumerate}
%  
%   \item Donnez une règle de distribution permettant d'affecter chaque document
%   de $D$ à un serveur. Indiquez, selon cette règle, où se trouvent les documents
%   $id=13$ et $id=16$. 
%   \item On ajoute un serveur dont la valeur de hachage vaut 9. Où se place-t-il
%   dans l'anneau, et quelles sont les opérations à effectuer à son arrivée\,?
%   \end{enumerate}
%   
% \begin{solution}
%  (1) On applique modulo(8) sur l'id, et on stocke l'objet sur le serveur qui
% précède sur l'anneau (ou qui suit, cela revient au même). (2) Le serveur se
% place en position 1. Tous les objets stockées entre 1 et 2, venant du serveur 7,
% doivent être placés sur le nouveau serveur (dans l'hypothèse ``précède'').
% \end{solution}
%    On se pose la question de l'emplacement de la table de hachage (donnant
%   l'équivalence entre les IP des serveurs et leur emplacement sur l'anneau).
%   
% \begin{enumerate}
%   \item  On
%   choisit initialement la solution suivante\,: chaque serveur connaît
%   l'emplacement de son successeur sur l'anneau. Quel est le nombre de messages
%   dans le pire des cas pour trouver l'emplacement d'un objet\,?
%   \item Vous arrivez et proposez une solution de type \textit{Chord}, comme vu
%   en cours. Quels sont les amis du serveur à l'emplacement 4\,?
% \end{enumerate}
% \begin{solution}
%  (1) complexité linéaire\,: il faut parcourir tout l'anneau dans le pire des
%  cas. (2) les amis sont 6 et 7. 
% \end{solution}
% 
\end{document}
